\documentclass[11pt]{article}
\usepackage[T1]{fontenc}
\usepackage{amsmath,amssymb,geometry}
\geometry{margin=1in}
\title{Rank zero has a degree-one identity\\Disproof of Conjecture 00000002766}
\author{ziangni-sys}
\date{October 8, 2026}
\begin{document}
\maketitle
\paragraph{Stated scope.}
The conjecture concerns polynomial identities of the variety of
matrices of rank at most $r$, and asserts that its minimal identity
degree is $2r+2$. Neither the English nor the Chinese version requires
$r>0$. We disprove this clause at the admissible rank bound $r=0$.
The coefficient-growth clause need not be resolved to disprove the
conjunction.

\paragraph{The actual matrix variety.}
Work over $\mathbb Q$. For every ambient size $n\geq1$,
\[
 V_0(n)=\{A\in M_n(\mathbb Q):\operatorname{rank}(A)\leq0\}
       =\{0\}.
\]
Indeed the rank is the dimension of the image of the linear map
$x\mapsto Ax$. Dimension zero forces that image to be the zero
subspace, hence the linear map vanishes. Evaluating it on every
standard basis vector shows that every column of $A$ is zero.
The converse is immediate. Thus this description covers every
matrix in the actual rank-constrained variety, at every ambient size.

\paragraph{A genuine polynomial identity.}
Take $p(X)=X\in\mathbb Q[X]$. This is a nonzero polynomial of degree
one. A one-variable polynomial is also a noncommutative polynomial
identity candidate: its evaluation in a matrix algebra is well-defined
by sending scalars to scalar matrices and $X$ to the input matrix.
For all $n$ and all $A\in V_0(n)$ its actual evaluation is
\[
 p(A)=A=0.
\]
Consequently a nonzero degree-one identity holds on the full rank-zero
matrix variety, uniformly in ambient size. Its minimum identity degree
is therefore at most one, whereas the asserted formula gives
$2(0)+2=2$. This is a contradiction.

\paragraph{Conclusion and formal coverage.}
The conjecture is false without an additional positive-rank hypothesis.
This proof makes no assertion about such a modified version.
The Lean project uses the actual \texttt{Matrix.rank}, defined by the
dimension of the matrix linear map's range. It proves rank at most zero
if and only if the matrix is zero, verifies the polynomial's nonzero
status and degree, and evaluates that polynomial through the actual
scalar-to-matrix algebra map at every square matrix size. Its final
theorem refutes the necessary degree lower bound that the proposed
minimum would entail. No assumed rank-zero certificate or scalar
replacement for the matrix variety is used.
\end{document}
